\section{Conclusion}

\textit{Written 2026-07-05.}

We presented a controlled study of how the choice of pretraining objective governs what transfers from an abundant vision-captured skeleton source (NTU RGB+D) to a data-scarce body-worn inertial target (Xsens), over a shared Local Relative Quaternion representation and a single Transformer encoder, under deployment-realistic subject-held-out few-shot calibration.

Four findings stand. \textbf{First}, objective utility is gap-gated: a 45 pp within-domain accuracy spread compresses to \textasciitilde{}4 pp across the modality gap, so the gap --- not the objective --- sets how much objective choice matters. \textbf{Second}, and against the field's working assumption that reconstruction/self-supervision is the safe default under shift, the significant negative-transfer case in this cross-modal setting is \textit{pure masked reconstruction}, which transfers below random initialization at every shot count (clip-level McNemar p<.001, seed-stable), while the hybrid reconstruction+supervision objective is robustly best and pure supervision is competitive and best-calibrated. \textbf{Third}, the mechanism is representational but two-part: a small feature-space gap is necessary but not sufficient (MMD does not track transfer), layer-wise CKA orders the objectives cleanly and monotonically with depth, yet alignment alone does not predict accuracy --- the objective's discriminative content decides transfer. \textbf{Fourth}, the prior's real deployment value is not headline accuracy but \textit{calibration and data-efficiency}: it calibrates better, converges 3--4 epochs faster, and delivers its largest benefit (+27 pp) precisely when only 0--1 subjects are enrolled, washing out by four.

Closing the human--machine-systems loop, an abstract controller driven by the recognizer's actual posterior stream shows these compressed accuracy differences compound sharply into task-level ones. A three-lock robustness protocol shows this compounding is not an artifact of one configuration, but also shows no objective is unconditionally safe under every stress test: under randomized task design the reconstruction objective compounds worst on average, while under a fixed, reliability-ranked task design the best-calibrated objective's own rare, confident misclassifications dominate once safety-critical errors are penalized severely --- calibration governs the safety/throughput frontier, but carries its own tail risk. An external skeleton$\rightarrow$skeleton check (CZU-MHAD) shows the negative-transfer signature is specific to the cross-modal gap, not the objective in the abstract.

The practical takeaway for skeleton-primed wearable recognition: prefer a hybrid objective, expect the prior to pay off mainly in calibration and low-enrollment efficiency rather than full-data accuracy, and --- for safety-sensitive control --- weight calibration quality alongside accuracy while accounting for confident-failure tail risk, not accuracy or calibration alone. Future work is the cross-modal external replication (the CZU body-worn IMU data), a full multi-seed contrastive sweep, and a calibrated, real-time controller instantiation.
